\section{Steps for Solving the Problem}

\subsection{Gaussian EDA as the Population Search}

The implementation uses the \texttt{GaussianUMDA} interface from \texttt{pateda}. Its search model is a product of univariate normal distributions. If $S_t$ denotes the selected individuals in generation $t$, the basic model can be expressed as
\begin{equation}
 p_t(\mathbf{x})=\prod_{j=1}^{d}\mathcal{N}(x_j;\mu_{t,j},\sigma_{t,j}^{2}),
 \quad \mu_{t,j}=\frac{1}{|S_t|}\sum_{\mathbf{x}\in S_t}x_j,
\end{equation}
where each variance describes the spread of the selected values in one coordinate. This factorization is computationally simple, but it cannot represent correlations between coordinates. That limitation is relevant on landscapes such as Rosenbrock, where good solutions follow a curved valley.

Each run uses a population of 30, selects the best half, and enables elitism. The initial population and subsequent generations are managed by the EDA framework within the specified bounds. Because the framework maximizes fitness, the minimization objective is supplied as $F(\mathbf{x})=-f(\mathbf{x})$. All reported results are converted back to the original objective, so smaller values always indicate better solutions. The plain method evaluates its initial population and 19 subsequent populations, consuming $30+19\times30=600$ evaluations.

\subsection{Shared Local-Search Interface}

The two hybrids retain the same population model, selection ratio, population size, and elitism setting. A local-search component receives a population and its evaluated fitness values, selects the member with the highest fitness, and attempts to improve it. The selected individual is replaced only if a strictly better evaluated candidate is found. Its original fitness is retained otherwise. This rule prevents a local call from directly worsening that individual, although the hybrid can still finish worse than the plain EDA because it allocates evaluations differently.

Both hybrids request 14 generations after initialization, using
\begin{equation}
 G_{\mathrm{hybrid}}=\left\lceil\frac{600}{30+10}\right\rceil-1=14.
\end{equation}
The local component is attached through the framework's internal component interface. Each invocation receives the remaining global budget and is capped at ten evaluations, both per generation and per selected solution. With only one selected solution, these two local limits coincide. The saved results confirm total use of 600 evaluations in every hybrid run. They do not separately record the number spent in the population and local components, so no measured breakdown between those components is claimed.

A counter wraps the true objective and is shared by the population and local searches. Every true call increments it, and a call after the global limit raises an error. Consequently, numerical differentiation and repeated evaluations count toward the same budget. Surrogate fitting and acquisition optimization do not call the benchmark objective and therefore do not consume this budget, although they consume computation time.

\subsection{Classical Refinement with L-BFGS-B}

The classical operator calls \texttt{scipy.optimize.minimize} with the L-BFGS-B method and the original global bounds. L-BFGS-B is a bound-constrained quasi-Newton method. No analytic Jacobian is supplied, so gradients are estimated numerically. The implementation sets both \texttt{maxiter} and \texttt{maxfun} to the available local budget. These parameters alone are insufficient to enforce a strict objective-call limit when numerical gradients are used, as documented by SciPy \cite{scipylbfgsb}.

An additional guard therefore returns a large penalty, $10^{100}$, after ten true local evaluations have been consumed. These subsequent penalty returns do not evaluate the benchmark. The wrapper tracks the best actual point visited and returns that point, rather than relying on the optimizer's final status or final iterate. This protects the evaluation budget but can interrupt a gradient estimate or line search. Accordingly, the experiment evaluates a tightly budgeted L-BFGS-B wrapper, rather than L-BFGS-B run to convergence.

In five dimensions, a forward-difference gradient can require approximately six objective calls including its base point, leaving little room for line searches. The classical wrapper charges repeated starting-point evaluations to its budget, whereas BO reuses the known centre.

\subsection{Bayesian Refinement with Expected Improvement}

BO constructs a probabilistic surrogate from evaluated points and uses an acquisition function to choose the next evaluation \cite{frazier2018tutorial}. Here each local call creates a new \texttt{skopt.Optimizer} with a Gaussian-process surrogate and Expected Improvement (EI). Gaussian processes provide predictive means and variances from a covariance model \cite{rasmussen2006gp}. The notebook leaves the detailed covariance and fitting settings at the library defaults; these are not independently tuned in this study.

The search region is a fixed box around the selected centre $\mathbf{c}$:
\begin{equation}
 a_j=\max\{\ell_j,c_j-0.05(u_j-\ell_j)\},\qquad
 b_j=\min\{u_j,c_j+0.05(u_j-\ell_j)\}.
\end{equation}
Thus, 5\% is the radius along each coordinate. Away from a global boundary, the full local width is 10\% of the global range. This is a fixed local box, not an adaptive trust-region algorithm.

The already evaluated centre and its objective value initialize the optimizer through \texttt{tell}; no extra benchmark call is needed. Setting \texttt{n\_initial\_points=1} makes this observation satisfy the initial-design requirement. Each of the next ten iterations asks for a candidate, evaluates its true objective, and updates the surrogate. For minimization, EI can be written as
\begin{equation}
 \mathrm{EI}(\mathbf{x})=\mathbb{E}\!\left[\max\{0,f_{\min}-f(\mathbf{x})-\xi\}\mid\mathcal{D}\right],
\end{equation}
where $\mathcal{D}$ contains the observations, $f_{\min}$ is the best observed local value, and $\xi$ is the acquisition's exploration parameter. The implementation uses the library's default acquisition parameters. A new surrogate is created at every local invocation, so previous generations' observations are not reused. Each completed call therefore has only the centre and at most ten newly evaluated observations available to guide refinement.

